%\documentclass[AMA,STIX2COL,Linenumberson]{MRM}
\documentclass[12pt,letterpaper,onecolumn]{article}

\usepackage{amsmath}
\usepackage{setspace}
\usepackage{graphicx}
\usepackage{amssymb}
\usepackage[superscript]{cite}
\usepackage{color}
\usepackage[colorlinks,citecolor=blue]{hyperref}
\usepackage{caption}
\usepackage{subcaption}
\usepackage{comment}
\usepackage[normalem]{ulem}

\usepackage{todonotes}
\setlength{\marginparwidth}{1.4cm}
\newcommand{\mycomment}[1]
{\todo[caption={#1},color=red!100!green!33,noline]{\color{black} \begin{spacing}{0.5}#1\vspace{-0.5em}\end{spacing}}}

%\renewcommand{\mycomment}[1]{}

\def\baselinestretch{1.5}
\oddsidemargin 0.0in
\textwidth 6.5in
\textheight 9.0in
\topmargin -0.4in
\headheight 0.0in

\makeatletter
  \renewcommand\@biblabel[1]{#1.}     % Arabic numbers, no brackets
\makeatother
\renewcommand{\refname}{\MakeUppercase{REFERENCES}}

\begin{document}


\begin{center}
{\LARGE\bf An Evaluation of Brain Volume and Cortical Thickness Measurement at 0.55T}
\end{center}
 
\vfill

\noindent Anand A. Joshi$^{1,3}$, Sophia X. Cui$^2$, Chin-Cheng Chan$^1$, Hao-Ting Kung$^1$, Jonas T. Kaplan$^3$, Krishna S. Nayak$^1$, Richard M. Leahy$^1$,  Justin P. Haldar$^{1,3}$

    
\vfill
\vspace{.1in}


\begin{singlespace}
    \noindent $^1$Signal and Image Processing Institute, Ming Hsieh Department of Electrical and Computer Engineering, University of Southern California, Los Angeles, CA, USA\\

	\noindent $^2$Siemens Medical Solutions USA, Los Angeles, CA, USA\\

	\noindent $^3$Brain and Creativity Institute, University of Southern California, Los Angeles, CA, USA
\end{singlespace}

\vfill 

\begin{singlespace}
	\noindent {\bf Address Correspondence to:}
	\vspace{.25in}
		
	\noindent Anand Joshi
			
	\noindent University of Southern California
	
	\noindent University Park Campus
	
	\noindent 3740 McClintock Ave,  EEB 426
	
	\noindent Los Angeles, CA 90089
	
	\noindent Tel: (213)740-4676
	
	\noindent Email: ajoshi@usc.edu
\end{singlespace}

\vfill

\noindent MANUSCRIPT BODY APPROXIMATE WORD COUNT:  3000

\newpage


\section*{ABSTRACT}

\noindent \textbf{Purpose:} To determine the feasibility of anatomical neuroimaging studies using 0.55T MRI.

\noindent \textbf{Methods:} MPRAGE images with 1 mm isotropic resolution were acquired from a cohort of five healthy subjects using both 0.55T and 3T scanners, including repeat scans of each subject at each field strength.  The 0.55T images were denoised to mitigate the effects of high resolution and low field strength.  All scans were analyzed using BrainSuite and FreeSurfer image analysis software, which performs operations including tissue classification, region segmentation/labeling, and thickness/volume measurements for different brain regions. Regional cortical thickness measurements and regional brain volume measurements were compared, and test-retest repeatability at each field strength was evaluated.

\noindent \textbf{Results:} For both image analysis software tools, we observed a 
 relatively strong correlation between field strengths for cortical thickness measures ($R^2 > 0.5$) and an excellent correlation between field strengths for regional brain volume measures ($R^2 > 0.99$).   Test-retest repeatability analysis suggests that experimental variability was generally higher at 0.55T than at 3T, although the amount of difference was dependent on the specific measure being considered and the software tool that was used.

\noindent \textbf{Conclusion:} 
 While our sample size was small which limits the strength of the conclusions we can draw, this {\color{.} small-scale pilot} study suggests the potential feasibility of certain types of anatomical neuroimaging studies at 0.55T (especially those involving volume measurements), assuming that appropriate adjustments are made to account for increased experimental variability at this field strength.


\section*{KEYWORDS}
{0.55T MRI; brain image analysis; cortical thickness; regional brain volume}

%TC:endignore


\clearpage

 
\section{INTRODUCTION}\label{sec1}
In recent years, the MRI community has increasingly scrutinized the limitations of modern higher-field MRI systems, particularly in terms of cost, accessibility, and compatibility, which pose significant challenges to widespread adoption and equitable healthcare access. This has led to renewed interest in lower field strengths \cite{sarracanie2015,campbellwashburn2019,marques2019,wald2020,marques2021,arnold2023,campbellwashburn2023,campbell2021feasibility}, with many contemporary low-field systems now available in the commercial marketplace.  However, while lower-field systems can have certain advantages, a key limitation is that the MR signal is typically weaker at lower fields, due to lower equilibrium polarization, which often results in reduced image quality compared to higher-field acquisitions.   Specifically, for a fixed experiment duration and depending on protocol design choices, lower-field acquisitions will often yield lower signal-to-noise ratio (SNR) or coarser spatial resolution than higher-field counterparts.

Modern MR neuroimaging studies frequently depend on software tools  that automate many aspects of brain image analysis, such  as FreeSurfer \cite{fischl2012freesurfer}, FSL \cite{jenkinson2012}, SPM \cite{ashburner2009}, ANTS \cite{tustison2021}, and BrainSuite \cite{shattuck2002brainsuite}.  Such software tools can perform various useful operations, including identifying and extracting the brain parenchyma from MR images of the head, classifying different brain tissue types (e.g., white matter, gray matter, and cerebrospinal fluid), co-registering the brains from different subjects into a common atlas space so that they can be more easily compared, subdividing the brain into labeled regions of interest, and quantifying the sizes of different anatomical features (e.g., measuring the thickness of the cortical gray matter in different regions or measuring the volume of different white matter structures).  Historically, these software tools have generally been developed and tuned to work with the type of high-quality human head images that are typically produced by higher-field MRI systems (e.g., $\sim$1 T or above).


Motivated by the potential opportunity to perform MR neuroimaging studies using modern lower-field MRI systems, the present {\color{.}small-scale pilot} study investigates the performance of brain image analysis software tools (specifically BrainSuite \cite{shattuck2002brainsuite} and FreeSurfer \cite{fischl2012freesurfer}) on brain images acquired with a contemporary 0.55T system, with results compared against those obtained from 3T images of the same subjects.  Our results provide evidence supporting the potential for using 0.55T MRI data in anatomical neuroimaging studies. \textcolor{.}{Specifically,  we examine the extent to which the cortical thickness and ROI volume measures are matched (or might be linearly transformed to match) across field strengths.}

Previous studies have also examined the performance of brain image analysis software with lower-field MRI data \cite{deoni2021,iglesias2023,cooper2024}.  We believe that our study is complementary to this growing body of work, including the exploration of new dimensions.  Specifically, we evaluated performance with 0.55T data, while previous studies focused on data from a much lower field strength (0.064T).  In addition, previous studies used a low-resolution data acquisition strategy for their low-field images (with some studies relying on super-resolution image enhancement techniques).  In contrast, we acquired noisy high-resolution data that was subsequently denoised using a regularized reconstruction technique that yields SNR improvements at the expense of spatial resolution.  

Our approach was motivated by the potentially surprising fact that denoising a noisy high-resolution image can have major efficiency advantages over a more conventional low-resolution high-SNR acquisition \cite{haldar2008, haldar2011a, haldar2011c, haldar2013} (see Ref.~\citen{haldar2022} for discussion of historical misconceptions about denoising versus super-resolution and the trade-off between resolution and SNR). Our work is also complementary to the substantial body of work evaluating the performance of brain image analysis software as a function of field strength, which has thus far focused on higher-field ($\geq$1.5 T) contexts \cite{dickerson2008,jovicich2009,lusebrink2013,heinen2016,isaacs2020,srinivasan2020}.

  
 
\section{METHODS}\label{sec:methods}
\subsection{Data Acquisition}
T1-weighted MPRAGE images with 1 mm isotropic resolution were collected from five healthy volunteers at 0.55T and 3T.  These images were acquired during the same scan sessions described in a previous study of 0.55T  diffusion MRI \cite{kung2024diffusion}.   The 0.55T datasets were acquired using a whole-body 0.55T system (prototype MAGNETOM Aera, Siemens Healthineers, Forchheim, Germany) equipped with high-performance shielded gradients (45 mT/m amplitude, 200 T/m/s slew rate), and the 3T datasets were acquired with a standard 3T system  (MAGNETOM Prisma$^{\text{fit}}$, Siemens Healthineers, Forchheim, Germany).  Each subject was scanned twice at each field strength, exiting the magnet between scans. All scans were performed under a protocol approved by our Institutional Review Board, after providing written informed consent.

The 3T scans were performed using a 20-channel receiver array, using one of our 3T center's standard neuroscience protocols  (TE=3.13~msec, TR=2530~msec, TI=963~msec, echo spacing=7.4~msec, bandwidth=220~Hz/Px, flip angle=10$^\circ$, matrix size=256$\times$256$\times$208, no parallel imaging acceleration, total acquisition time=10~min~48~sec). \textcolor{.}{Note that while this acquisition time is longer than that in many popular protocols (e.g., those from ADNI or the Human Connectome Project), this is the default protocol at our 3T imaging center and is designed to produce high-quality images in cooperative subjects where motion artifacts are minimal.}

The 0.55T scans were acquired with a 16-channel head and neck receiver array coil.  We used different sequence parameters at 0.55T than we did at 3T because of contrast differences (relaxation parameters change with field strength) and hardware differences (the 3T system had a more powerful gradient system, with 80 mT/m amplitude and 200 T/m/s slew rate).  The 0.55T sequence parameters we employed in this study (TE=3.43~msec, TR=1700~msec, TI=858~msec, echo spacing=8~msec, bandwidth=190 Hz/Px, flip angle = 20$^\circ$, 1 mm isotropic voxel size, no parallel imaging acceleration) were selected heuristically based on data from pilot scans, where we used a subjective visual assessment of gray-white tissue contrast to select the best option from a small number of (non-exhaustive) protocol variations.    Since the 0.55T scan had a shorter TR than the 3T scan, we used 50\% phase oversampling to match the total acquisition time of the 3T scan (10~min~48~sec).  The use of phase oversampling is equivalent to using a larger FOV along the phase encoding dimension and achieves the same SNR benefits as traditional data-averaging but with finer control over the total acquisition time.  


\subsection{Image Reconstruction}

As illustrated in Fig.~\ref{fig:attn}(a), the 0.55T acquisition protocol we chose yields low SNR images if only standard Fourier image reconstruction is employed. {For example, standard Fourier reconstructions from one representative subject were observed to have average white-matter SNRs of approximately 6 at 0.55T, which was substantially lower than the corresponding SNRs of approximately 30 observed at 3T. Note that at both field strengths, SNR varied spatially because of the spatially-varying sensitivities of the receiver coils -- tending to be lower than these average values near the center of the brain and higher near the cortex.   }

This level of SNR is expected to be problematic for automated brain image analysis since most analysis software (including BrainSuite and FreeSurfer) are designed for higher-quality input images.\footnote{\color{.}Indeed, applying these image analysis pipelines directly to the noisy images produces substantially suboptimal segmentation and surface extraction results that are too poor to warrant further analysis, as should be expected. Interestingly, we also observed that our attempts to apply SynthSeg or SynthSR \cite{billot2023synthseg, iglesias2023synthsr} (AI-based approaches that are distributed with FreeSurfer and are designed to segment arbitrary brain MRI scans or synthesize high-quality T1 scans from arbitrary lower-quality brain MRI scans) were unsuccessful with this data, and produced results that were worse than directly using the noisy images. We suspect that this is because we used an atypical high-resolution low-SNR acquisition scheme these techniques were nor designed for. } Our choice to intentionally acquire noisy high-resolution data may also be viewed as somewhat unorthodox since conventional wisdom would prefer a higher-SNR lower-resolution acquisition\cite{haldar2022}.  However, as mentioned in the introduction,  theoretical results suggest that major SNR/resolution efficiency benefits can be obtained by applying denoising reconstruction to low-SNR high-resolution data \cite{haldar2008, haldar2011a, haldar2011c, haldar2013}. {\color{.} Our 0.55T acquisition was designed with this type of approach in mind, and as such, we implemented an image reconstruction approach that was designed to enhance the SNR of our high-resolution 0.55T data.} {\color{.} Note that this reconstruction approach was only applied to the 0.55T data and not the 3T data, since the 3T images did not need SNR enhancement and might be negatively impacted by the loss of spatial resolution associated with regularized reconstruction. }
 \textcolor{blue}{Note that this reconstruction approach was only applied to the 0.55T data and not the 3T data, since the 3T images did not need SNR enhancement. While applying the regularized reconstruction to the 3T data might present an interesting methodological comparison to demonstrate the impact of denoising on high-field data, we chose to focus our evaluation on the standard high-quality 3T acquisitions representing the current clinical and research standard. This choice avoids unnecessarily compromising the spatial resolution of the 3T data due to the smoothing effects of regularization.\mycomment{R1.1}}

 
Our in-house reconstruction implementation employed two complementary SNR enhancement approaches.  We first used coil compression to remove uninformative noise-dominated virtual channels from the original multichannel k-space data \cite{buehrer2007,huang2008,kim2021} (we specifically used the ROI-SVD technique from Ref.~\citen{kim2021} to focus the virtual channels on spatial regions containing brain parenchyma, selecting the number of virtual channels to preserve at least 98\% of the signal energy from these regions).  We then performed regularized reconstruction from the coil-compressed k-space data, using a SENSE-based data acquisition model\cite{pruessmann2001} (with sensitivity maps estimated using PISCO \cite{lobos2023}) together with (smoothing-based) total-variation regularization\cite{block2007}.  

An illustrative example of regularized reconstruction is shown in Fig.~\ref{fig:attn}(b).   As can be seen, the appearance of noise is substantially reduced in the regularized reconstructions. However, we did not attempt to quantify the SNR improvement because of the nonlinear nature of total variation regularization.  

After reconstruction, the images were bias-field corrected using N4 \cite{tustison2010} and corrected for gradient nonlinearity using software from the Human Connectome Project's processing pipeline \cite{glasser2013minimal}.

\subsection{Image Analysis using BrainSuite and FreeSurfer}
We processed all of the 0.55T and 3T images from all subjects using the anatomical image analysis pipelines from BrainSuite \cite{shattuck2002brainsuite} (Version 23a) and  FreeSurfer \cite{fischl2012freesurfer} (Version 7.4.1).  In both cases, the anatomical pipeline takes an input T1-weighted MRI volume and performs a series of processing steps to generate cortical mesh models (skull-stripping, bias field correction, tissue classification, inner cortical mesh generation, and expansion of this mesh to generate a pial cortical mesh).  Subsequently, the brain volume and cortical surface are segmented by coregistration to a pre-labeled anatomical atlas.  Illustrative cortical surfaces and labeled surfaces and volumes produced by this process are shown in Fig.~\ref{fig:output_brainsuite}.

The thicknesses of different cortical regions of interest (ROIs) and the volumes of different volumetric brain ROIs are calculated as endpoints of this process.  We choose to focus our performance evaluation on these thickness and volume measures, as these endpoints are influenced by all aspects of the anatomical processing pipeline and are also quantities of interest in many neuroscience studies.

\subsection{Performance Evaluation}
We performed several statistical analyses to assess the degree of consistency between the 0.55T results and the 3T results and to assess the degree of repeatability for thickness and volume measures at each field strength.

First, we used simple linear regression models to understand the degree to which the thickness and volume measurements at 0.55T could be predicted by the corresponding 3T thickness and volume measures from the same ROIs in the same subject, pooling all volumetric ROIs and subjects together into a single analysis for the volume prediction, and pooling all surface ROIs and subjects together into a single analysis for the thickness prediction. For these analyses, the measures from repeat scans were averaged. The goodness of fit was assessed using the coefficient of determination ($R^2$).  

{Note that linear regression assesses the correspondence between two measurements while accounting for potential systematic linear differences between them, and can be  appropriate in scenarios (such as this one) where systematic differences may be expected.  However, it can also be informative to evaluate agreement without accounting for systematic differences.  To that end, we also performed Bland-Altman analyses\cite{altman1983measurement} comparing 0.55T thickness and volume measurements against the corresponding 3T measurements.} This included a calculation of the bias (the mean of the difference) between the two measurements, as well as the 95\% limits of agreement (the interval within which 95\% of the data points are expected to fall). \textcolor{blue}{The limits of agreement were calculated as the mean difference $\pm 1.96$ times the sample standard deviation (SD) of the differences, rather than as confidence intervals for the mean difference.\mycomment{R1.3}} 

We also used Bland-Altman plots  to evaluate the test-retest repeatability of thickness and volume measures at both 0.55T and 3T. \textcolor{blue}{For the ROI volume Bland-Altman plots, we additionally identified outlier regions using a predefined relative-deviation criterion. For each ROI and subject, the absolute Bland-Altman difference was divided by the corresponding Bland-Altman mean volume. An ROI was considered an outlier region if this relative deviation, averaged across the five subjects, exceeded 10\%. In the plots, all subject-level points belonging to these outlier ROIs were colored orange; region names are reported in the Results rather than labeled directly in the panels to preserve readability.\mycomment{R1.3}}

\section{RESULTS}
\label{sec:results}
The results of the linear regression analyses are shown in Fig.~\ref{fig:scatter_plots}. As can be seen, for both BrainSuite and FreeSurfer, the thickness and volume measures obtained at 0.55T were { relatively well-predicted} by the corresponding 3T values, with especially good consistency for volume measures ($R^2 > 0.50$ for thickness measures and $R^2>0.99$ for volume measures in both cases).  Note that the BrainSuite and FreeSurfer measures are not directly comparable since they use different atlases that each segment the brain and the cortical surface in distinct ways (i.e., BrainSuite uses the BCI-DNI atlas, with 33 cortical ROIs per hemisphere and 29 noncortical ROIs \cite{joshi_hybrid_2022}, while FreeSurfer  uses the Desikan-Killiani atlas, with 34 cortical ROIs per hemisphere and 14 noncortical ROIs \cite{desikan_automated_2006}). 

The results of the 0.55T versus 3T Bland-Altman analyses are shown in {\color{.}Fig.~\ref{fig:bland-altman-p55T-vs-3T}}. As can be seen, there was relatively good agreement between 0.55T and 3T, although small constant/linear bias trends are potentially present for some measures (which were modeled in the previous regression analyses but are treated as discrepancies in Bland-Altman analysis, which does not account for systematic differences). \textcolor{blue}{For the ROI volume panels, the x-axis is shown on a logarithmic scale to improve visualization across the wide range of regional volumes. Orange points in panels (c,d) indicate subject-level measurements from outlier ROIs, defined as ROIs for which the average across-subject absolute relative field-strength volume difference, $|\Delta V|/\bar{V}$, exceeded 10\%. These outlier regions included orbitofrontal/subcallosal/temporal-pole regions, claustrum, lateral geniculate nucleus, superior colliculus, mammillary body, ventricular/choroid-plexus/vessel labels, cerebellar white matter, accumbens, optic chiasm, and selected cortical regions. These outlier names are summarized here rather than labeled directly in the panels to preserve figure readability.\mycomment{R1.3}}



The test-retest Bland-Altman plots for ROI volume measures are shown in Fig.~\ref{fig:bland-altman-roi_vols}.  As can be seen, the bias was close to zero (\textcolor{blue}{$< 0.07$ cm$^3$\mycomment{R1.3}}) in all cases.    The 95\% limits of agreement intervals were roughly 50\% wider at 0.55T than they were at 3T for BrainSuite, and roughly 70\% wider at 0.55T than they were at 3T for FreeSurfer.  This suggests that repeatability is worse at 0.55T than it is at 3T, as should be expected. \textcolor{blue}{The x-axis is plotted on a logarithmic scale. Orange points indicate subject-level measurements from outlier ROIs, defined as ROIs for which the average across-subject absolute relative test-retest difference, $|\Delta V|/\bar{V}$, exceeded 10\%. The corresponding outlier regions are summarized in the text rather than labeled directly in the panels to preserve readability.\mycomment{R1.3}} 

{\color{.}Notably, this test-retest analysis of ROI volume measures reveals some large outliers with intra-session differences greater than 5 cm$^3$.  Upon closer inspection, we found that these outliers occur almost exclusively in cerebellar regions.  Further, we also observed that total overall cerebellar volumes remained relatively stable between repetitions, and that the large deviations were driven by a misclassification and swapping of tissue labels between Cerebellar White Matter and Cerebellar Cortex. We suspect that this occurred because the cerebellum contains finely interleaved layers of gray and white matter (folia), and the lower data quality at 0.55T makes it difficult for automated segmentation algorithms to reliably distinguish these tissue borders.} \textcolor{blue}{These regions represent anatomical or methodological constraints where reliable measurement may not be achievable, independent of field strength.\mycomment{R1.3}} \textcolor{blue}{More generally, the 10\% outlier analysis identified several small or anatomically challenging regions, including small nuclei, vessels, choroid plexus/ventricular structures, temporal pole/entorhinal regions, and cerebellar white matter. The fact that some of these outliers are also present in the 3T test-retest plots suggests that they reflect anatomical or methodological constraints of automated regional labeling, rather than a limitation that is specific to low-field imaging.\mycomment{R1.3}}

The test-retest Bland-Altman plots corresponding to cortical thickness measures are shown in Fig.~\ref{fig:bland-altman-thickness}. Similar to the volume measures, the bias for cortical thickness was close to zero ($< 0.05$ mm) in all cases.   The 95\% limits of agreement intervals were roughly 20\% wider at 0.55T than they were at 3T for BrainSuite, and roughly 80\% wider at 0.55T than they were at 3T for FreeSurfer.  As before (and as expected), this suggests that repeatability is worse at 0.55T than it is at 3T. 

It is perhaps surprising that we see a bigger difference between 0.55T and 3T for the repeatability of cortical thickness (corresponding to a relatively small brain feature) than we do for the repeatability of ROI volumes (corresponding to a larger brain feature).  One possible explanation for this is that cortical thickness and ROI volume measurements have different levels of sensitivity to different aspects of the processing pipeline, with cortical thickness a bit more sensitive to the tissue classification step and ROI volumes a bit more sensitive to volume labeling/segmentation.   

\section{DISCUSSION AND CONCLUSIONS}\label{sec:conclusion}
This {\color{.} small-scale pilot} study provides insight into the feasibility of using high-resolution 0.55T MRI data for anatomical neuroimaging research, specifically for tasks involving cortical thickness and ROI volume measurements. {\color{.} This  adds to the modern literature on lower-field MRI, which has been receiving renewed interest as the MRI community reflects on some of the limitations of higher-field systems (such as cost, accessibility, and compatibility).}  We observed relatively strong correlations between thickness and volume measurements obtained at 0.55T and 3T. 
In addition, as might be expected, we observed slightly wider \textcolor{blue}{95\% limits of agreement (mean difference $\pm 1.96$ SD of the differences)\mycomment{R1.3}} for measures at 0.55T than we did at 3T.  Overall, although our sample size was relatively small, we believe that our results suggest that noisy high-resolution 0.55T MRI data {has the potential} to achieve sufficient quality to be used for {certain types of} anatomical neuroimaging research, {especially those involving volume measurements.}

\textcolor{blue}{Based on the statistical assumptions of a t-test, the number of subjects required to achieve significance is influenced by population variance, which includes both biological variability and experimental variability (including the effects of noise). Our results suggest that experimental variability is higher at 0.55T than at 3T, and this variance increase is proportionally larger for volumetric measures (~50\% to ~300\%) than for cortical thickness (25\% to 60\%). Consequently, if population variance is dominated by measurement noise, volumetric studies at 0.55T would require disproportionately larger cohorts to maintain comparable statistical power relative to 3T, whereas cortical thickness measurements may remain more feasible. However, although our sample size was small, we observed that true inter-subject biological variability at 3T was often substantially larger (e.g., often 10$\times$ or more in many ROIs) than test-retest variability at 0.55T. Because biological variance is typically the dominant factor, we hypothesize that the increased measurement noise at 0.55T may only have a modest impact on the overall population variance. Thus, comparable levels of statistical significance might be achievable at 0.55T with only a small increase in sample size relative to 3T. A larger follow-up study focusing on a specific scientific question is needed to definitively assess statistical power and sample size differences.\mycomment{R1.2}}

{\color{.} Our  results are based on using a high-resolution low-SNR acquisition, which is atypical for traditional MRI (where acquisitions are frequently designed to prioritize SNR), although is aligned with theoretical results demonstrating that low-SNR acquisitions can be surprisingly efficient when combined with appropriate denoising \cite{haldar2008, haldar2011a, haldar2011c, haldar2013}.  The results shown in this paper provide further empirical evidence that high-resolution low-SNR acquisitions can be more efficient than might be traditionally expected.} {\color{.} However, the experiments in this work relied on a fixed high-resolution acquisition protocol, which had not been directly optimized for efficiency.  Recent work suggests that experimental optimization can be very beneficial in scenarios like this \cite{wang2025kspace}, and a more detailed experimentally-driven investigation of alternative acquisition strategies would be an interesting topic for future work. }

For simplicity, our results were obtained using a very simple regularized reconstruction technique to denoise the data, representing a standard/widespread approach from the regularized image reconstruction literature.  Exploring other image enhancement strategies is beyond the scope of this paper, but the results observed using this simple approach are already encouraging, and we expect that more advanced approaches will be able to achieve even better results.

To keep the study practically manageable, we only evaluated the performance of two different brain image analysis pipelines (i.e., those of BrainSuite and FreeSurfer).  The results we obtained with BrainSuite and FreeSurfer were relatively consistent, and we hypothesize that similar results would have also been obtained from other analysis pipelines.

The images we acquired at 0.55T and 3T had {systematic differences in both contrast (due to the natural variation of relaxation characteristics with field strength) and perceived sharpness (due to the use of smoothness-based regularization). Image contrast and sharpness are very important factors in automated brain image analysis, and these systematic image differences  may contribute to systematic differences in tissue classification (e.g., the boundary between white and gray matter) and surface identification, which in turn may contribute to systematic differences in labeling, cortical thickness, and ROI volume measurements between 0.55T and 3T. Notably, the imaging parameters we used at 0.55T were chosen qualitatively and may not be optimal for this task, while we also used a relatively simple denoising reconstruction procedure that may not preserve sharp edges as well as more advanced noise-reduction approaches. It may be of interest for future studies to more carefully examine the impact of different contrast-related acquisition parameters and different noise mitigation strategies.}  In addition, while systematic differences in thickness and volume estimates are not unexpected, our linear regression analysis demonstrates that these differences can {potentially} be calibrated fairly well with a regression line {(i.e., we observed $R^2>0.5$ in all cases with $R^2>0.7$ in three out of the four cases)}, which could potentially allow for the approximate ``conversion" of thickness and volume measurements from one field strength to another. { \textcolor{.}{However,  additional work would be required to validate this hypothesis before we would recommend deploying such an approach in a practical application, particularly given that the variability of cortical thickness and ROI volume measures may be different for different brain structures.} In addition, more advanced harmonization techniques could be an interesting topic for future exploration.}


\section*{Data Availability}  
Data and source code for the study will be shared, where possible,  upon reasonable request.  

\section*{Declarations}
\textbf{Conflict of interest:} Sophia X. Cui is an employee of Siemens Medical Solutions USA, and this work received research support from Siemens Healthineers.
All other authors declare that they have no conflicts of interest. 

\noindent \textbf{Ethics Approval:} This study was performed with approval from the Institutional Review Board of the University of Southern California. Written informed consent was obtained from all participants.


\section*{ACKNOWLEDGMENTS}
This work was supported in part by NIH grant R56-EB034349, NSF grant CNS-1828736, DoD grant HT94252310149, the Ming Hsieh Institute for Research on Engineering-Medicine for Cancer, computing resources from the USC Center for Advanced Research Computing, and research support from  Siemens Healthineers.


\begin{comment}
\section*{CONFLICT OF INTEREST STATEMENT}
Sophia X. Cui is an employee of Siemens Medical Solutions USA, and this work received research support from Siemens Healthineers.

\end{comment}

\bibliographystyle{MRMbib}
\bibliography{references}

\clearpage

\section*{FIGURES}



\begin{figure}[ht]
  \centering
  \begin{subfigure}[b]{0.24\columnwidth}
      \centering
      \includegraphics[width=1\textwidth]{original.png}
      \caption{}
   \end{subfigure}
   \begin{subfigure}[b]{0.24\columnwidth}
      \centering
      \includegraphics[width=1\textwidth]{denoised.png}
      \caption{}
    \end{subfigure}
    \caption{Illustrative examples of (a) conventional Fourier reconstruction (with root sum-of-squares coil combination) of the 0.55T data and (b) total-variation regularized reconstruction of the same data.}
    \label{fig:attn}
\end{figure}

\clearpage

\begin{figure}[t]
    \centering
    \includegraphics[width=\textwidth]{imgs_segs_v2.pdf}
    \caption{Example MRI scans at (a) 0.55T and (b) 3T, and the corresponding (c,d) volume segmentations, (e,f) mid-cortical surface reconstructions, and (g,h) cortical surface segmentations obtained using BrainSuite.
}
    \label{fig:output_brainsuite}
\end{figure}

\clearpage

\begin{figure}[t]
    \centering
    \includegraphics[width=0.95\linewidth]{scatter_plots_lin_regression_v4.pdf}
    \caption{Results of linear regression between the anatomical measures obtained at 0.55T and 3T.  We show (top row) cortical thickness and (bottom row) ROI volume results corresponding to (left column) BrainSuite and (right column) FreeSurfer processing.  In addition to showing scatter plots of the individual data points, we also show the best fit regression lines and the corresponding $R^2$ values.   
    }
    \label{fig:scatter_plots}
\end{figure}

\clearpage

\begin{figure}[hb]
    \centering
    \includegraphics[width=1\linewidth]{Bland_altman_3Tvs0.55T.pdf}
    \caption{Bland-Altman plots comparing 0.55 T measurements to the corresponding 3T measurements for (a,b) cortical thickness and (c,d) ROI volume measurements corresponding to (a,c) BrainSuite and (b,d) \textcolor{blue}{FreeSurfer} processing. \textcolor{blue}{Insets report the bias (mean difference) $\pm$ SD, together with the 95\% limits of agreement calculated as bias $\pm 1.96$ SD. The dashed red line denotes the bias and the dashed black lines denote the limits of agreement. The ROI volume panels (c,d) use a logarithmic x-axis. Orange points in panels (c,d) denote all subject-level measurements from outlier ROIs, defined as ROIs where the average across-subject absolute relative field-strength volume difference, $|\Delta V|/\bar{V}$, exceeded 10\%. The corresponding region names are summarized in the Results rather than labeled directly in the panels to preserve readability.}}
    \label{fig:bland-altman-p55T-vs-3T}
\end{figure}

\clearpage

\begin{figure}[t]
    \centering
    \includegraphics[width=1\linewidth]{bland_altmann_roi_vol_v7.pdf}
    \caption{Bland-Altman plots for ROI volumes corresponding to (a,b) BrainSuite and (c,d) FreeSurfer processing using  (a,c) 0.55T data and (b,d) 3T data. \textcolor{blue}{Insets report the bias (mean difference) $\pm$ SD, together with the 95\% limits of agreement calculated as bias $\pm 1.96$ SD. The dashed red line denotes the bias and the dashed black lines denote the limits of agreement. The x-axes use logarithmic scaling. Orange points denote all subject-level measurements from outlier ROIs, defined as ROIs where the average across-subject absolute relative test-retest difference, $|\Delta V|/\bar{V}$, exceeded 10\%. The corresponding region names are summarized in the Results rather than labeled directly in the panels to preserve readability.}}
    \label{fig:bland-altman-roi_vols}
\end{figure}

\clearpage

\begin{figure}[t]
    \centering
    \includegraphics[width=1\linewidth]{bland_altmann_thickness_v5.pdf}
    \caption{Bland-Altman plots for cortical thicknesses corresponding to (a,b) BrainSuite and (c,d) FreeSurfer processing using  (a,c) 0.55T data and (b,d) 3T data. \textcolor{blue}{Insets report the bias (mean difference) $\pm$ SD, together with the 95\% limits of agreement calculated as bias $\pm 1.96$ SD. The dashed red line denotes the bias and the dashed black lines denote the limits of agreement.} }
    \label{fig:bland-altman-thickness}
\end{figure}



\end{document}
